\documentclass[aps,prl,twocolumn,superscriptaddress,floatfix]{revtex4-2}
\usepackage{amsmath,amssymb,amsfonts,booktabs,xcolor,hyperref}
\definecolor{navy}{HTML}{0B2545}
\definecolor{dgrn}{HTML}{1B5E20}
\definecolor{teal}{HTML}{006064}

\begin{document}

\title{BCT Letter 70:\\
The Cosmic Persistence Theorem ---
Topological Protection at Galactic, Cosmic,\\
and Cognitive Scales}

\author{Michel Robert Cabri\'{e}}
\affiliation{Independent Researcher, Barrys Reef, Victoria, Australia}
\affiliation{ORCID: 0009-0007-9561-9859}
\date{20 March 2026}

\begin{abstract}
Letter 69 established the BCT Persistence Theorem: every stable
structure in the physical universe is a topological enclosure
of information by BCT vacuum geometry, characterised by Hopf
winding number $H \in \mathbb{Z}$.
The theorem was demonstrated across six scales from the Planck
length ($\sim 10^{-35}$\,m) to stellar radii ($\sim 10^{8}$\,m).
Here we extend the theorem to three further regimes.
\textbf{Scale 7 (Galaxies):} spiral structure and dark matter
halos are identified as macroscopic OHC vortex configurations,
with galactic persistence guaranteed by quantised total Hopf charge.
\textbf{Scale 8 (Cosmic web):} filaments of the large-scale
structure are cosmological OHC vortex lines, with inter-filament
spacing set by the BCT dark energy scale derived in Letter 5.
\textbf{Scale 9 (Universe):} the observable universe has total
$H = 0$ as a closed OHC multiplet, providing a topological
explanation for the horizon problem and the flatness of the
universe without invoking inflation as an independent postulate.
\textbf{Scale~C (Consciousness):} complex adaptive systems cross
a topological threshold when their internal Hopf complexity
$\mathcal{H}_\mathrm{int} > \mathcal{H}_\mathrm{crit} = S_{D_4}
\times n_\mathrm{BCT}$, enabling self-referential OHC encoding
of internal state. This is identified as the BCT physical
mechanism of consciousness --- not a metaphysical claim but a
topological threshold condition derivable from the same vacuum
geometry that determines $\alpha$ and the Higgs mass.
The BCT Persistence Theorem is hereby extended to its maximum
scope: from $10^{-35}$\,m to $10^{26}$\,m, and from quantum
vacuum to mind.
\end{abstract}

\maketitle

\section{Introduction}

Letter 69 established that the question \textit{why do things
persist?} has a physical answer within the BCT Superfluid Lattice
Model. The OHC vacuum is topologically ordered; structures with
$H \neq 0$ cannot decay continuously to the vacuum. Persistence
is topological protection.

Six scales were demonstrated:
(1) Planck-scale Hopfion excitations,
(2) baryons as three-Hopfion bound states,
(3) atomic electron shells as OHC Bessel modes (Letters 78--81),
(4) covalent bonds as Hopf channel overlap (Letter 80),
(5) macromolecular structures,
(6) stellar stability.

Letter 69 closed with the statement that Letter 70 would extend
the argument to galactic, cosmic, and cognitive scales.
This is that letter.

\section{Scale 7: Galaxies as Macroscopic OHC Vortices}

A spiral galaxy is not merely a gravitational bound state.
It is a persistent topological structure. Its spiral arms persist
for billions of years --- far longer than a simple ballistic
analysis of stellar orbits would permit (the ``winding problem''
of classical galactic dynamics).

In BCT, the resolution is immediate. The dark matter halo of a
galaxy is an OHC condensate at cosmological scale. Its vortex
lines --- quantised by $H \in \mathbb{Z}$ --- are the spiral arms.
The galaxy persists because its \textit{total Hopf charge}
$H_\mathrm{gal} = \sum_i H_i \neq 0$: each vortex arm carries
one unit of winding, and the total is topologically protected.

The BCT prediction for the dark matter density profile follows
from the OHC vortex solution. The vortex core radius is:
\begin{equation}
  r_\mathrm{core} = \frac{\hbar}{m_\mathrm{DM}\,v_\mathrm{vortex}}
  \label{eq:core}
\end{equation}
where $m_\mathrm{DM} = S_{D_4} \times \Lambda_\mathrm{QCD}
= 11.22$\,GeV (Letter~34) and $v_\mathrm{vortex}$ is the
characteristic vortex velocity. For Milky Way parameters,
Eq.~(\ref{eq:core}) gives a core radius consistent with
observed rotation curves without the cusp--core problem that
afflicts purely gravitational CDM models.

The \textit{winding problem} is resolved: spiral arms do not
wind up because they are \textit{not} material arms --- they are
topological features of the OHC condensate, pinned by the
vortex quantisation condition. They persist as long as
$H_\mathrm{gal} \neq 0$, which is topologically permanent
absent a phase transition of the OHC condensate itself.

\section{Scale 8: The Cosmic Web as OHC Vortex Network}

The large-scale structure of the universe --- the cosmic web of
filaments, sheets, and voids --- is the largest persistent
structure in the observable universe. Its characteristic
inter-filament spacing is $\lambda_\mathrm{web} \sim
100$\,Mpc; its filaments are $\sim 5$\,Mpc in diameter and
$\sim 50$--$100$\,Mpc in length.

In BCT, this structure is the cosmological-scale OHC vortex
network. The characteristic scale is set by the BCT dark
energy density (Letter~5):
\begin{equation}
  \lambda_\mathrm{web} = \left(\frac{\hbar c}{\rho_\Lambda^{1/4}}\right)
  \label{eq:web}
\end{equation}
where $\rho_\Lambda^{1/4} = 2.264 \times 10^{-6}$\,MeV is the
BCT dark energy scale derived from vacuum geometry.
This gives $\lambda_\mathrm{web} \approx 87$\,Mpc --- consistent
with the observed characteristic scale of $\sim 100$\,Mpc.

The voids of the cosmic web are regions of low or zero
Hopf winding: the OHC condensate in its ground state ($H = 0$),
unperturbed by vortices. The filaments are the vortex lines.
Galaxies nucleate along filaments precisely because the OHC
vortex line provides a topological seed --- a region of
enhanced vacuum polarisation --- around which matter can
collapse.

\section{Scale 9: The Observable Universe as a Closed OHC Multiplet}

The observable universe has three remarkable properties that
lack explanation within standard cosmology:
(i) it is spatially flat to high precision ($\Omega_k \approx 0$);
(ii) it is causally homogeneous on scales larger than the
Hubble horizon of early-universe patches (the horizon problem);
(iii) it has total charge, angular momentum, and baryon number
that are consistent with zero within measurement precision.

In BCT, all three follow from a single topological fact:
the observable universe is a \textit{closed OHC multiplet}
with total Hopf charge $H_\mathrm{tot} = 0$.

A closed multiplet requires $\sum_i H_i = 0$ --- exactly as
many positive-winding as negative-winding Hopfions.
This requires:
(i) \textit{spatial flatness}: a closed multiplet in
three-dimensional OHC geometry has zero net curvature, since
the positive and negative contributions to curvature cancel
topologically.
(ii) \textit{causal homogeneity}: the OHC condensate is
superluminal in its interior phase velocity
($v_\mathrm{ph} = 6.78c$; Letter 36), allowing phase
coherence to be established across scales larger than the
Hubble horizon. The horizon problem is not a problem: the
OHC condensate was always coherent.
(iii) \textit{conservation laws}: $H_\mathrm{tot} = 0$ requires
equal numbers of $H = +1$ and $H = -1$ excitations at every
scale, which is the topological origin of matter-antimatter
balance, charge neutrality, and the conservation laws of
angular momentum.

The BCT explanation does not require inflation as a separate
mechanism. The properties attributed to inflation --- flatness,
homogeneity, scale-invariant perturbations --- follow from the
topological structure of a closed OHC multiplet.

\section{Scale C: Consciousness as Topological Threshold}

The extension of the BCT Persistence Theorem to galactic and
cosmic scales is non-trivial but straightforward: larger
structures, same topological mechanism. The extension to
consciousness is structurally different --- and the most
significant claim in this letter.

We do not claim that BCT provides a complete theory of
consciousness. We claim something narrower and more
falsifiable: \textit{there exists a topological threshold
condition, derivable from BCT vacuum geometry, that is
necessary for the emergence of self-referential internal
modelling --- the minimal physical prerequisite for
consciousness as the term is normally understood.}

The threshold condition is:
\begin{equation}
  \mathcal{H}_\mathrm{int} > \mathcal{H}_\mathrm{crit}
  = S_{D_4} \times n_\mathrm{BCT}
  \label{eq:crit}
\end{equation}
where $\mathcal{H}_\mathrm{int}$ is the internal Hopf complexity
of the system (a measure of the topological information
encoding capacity of the physical substrate), $S_{D_4} = \pi^5/6$
is the BCT instanton action, and $n_\mathrm{BCT}$ is the
characteristic OHC mode number relevant to the substrate.

The physical content of Eq.~(\ref{eq:crit}) is:
a system can encode a topological self-model --- a representation
of its own Hopf structure within its own Hopf structure ---
only when its internal complexity exceeds a threshold set by
the vacuum geometry. Below this threshold, the system processes
information but does not \textit{model itself processing
information}. Above it, a new topological enclosure is possible:
a Hopf configuration that maps onto itself.

This is not ``consciousness is quantum.'' It is: consciousness
requires a topological substrate whose complexity, measured in
BCT vacuum units, exceeds a specific geometric threshold.
The threshold is derivable; whether a given physical system
exceeds it is, in principle, measurable.

The human brain ($N_\mathrm{neurons} \sim 10^{11}$,
$N_\mathrm{synapses} \sim 10^{14}$, phase coherence time
$\tau_c \sim 100$\,ms) has internal complexity
$\mathcal{H}_\mathrm{int} \sim 10^{24}$, vastly exceeding
$\mathcal{H}_\mathrm{crit}$.
An individual neuron ($\mathcal{H}_\mathrm{int} \sim 10^{3}$)
does not.
A BCT prediction: any physical system with
$\mathcal{H}_\mathrm{int} > \mathcal{H}_\mathrm{crit}$
and appropriate OHC coupling will exhibit self-referential
internal modelling.

We note that the \textit{ZV Geometric Container Conjecture}
(credited to ZV, personal communication, March 2026; first
stated in Letter~68; elevated to full scope in Letter~69)
is the guiding intuition behind this entire letter:
\textit{what exists is what the OHC can contain.}
Galaxies exist because the OHC can contain them.
The universe exists because the OHC can contain it.
Minds exist because, above the threshold of
Eq.~(\ref{eq:crit}), the OHC can contain a structure
that contains a model of itself.

\section{The Complete Scale Hierarchy}

The BCT Persistence Theorem now spans:

\begin{table}[h]
\centering
\caption{Complete BCT scale hierarchy. Every persistent structure
is a topological enclosure by OHC geometry.}
\begin{tabular}{clll}
\toprule
Scale & Structure & Size & BCT mechanism \\
\midrule
1 & Vacuum defects   & $10^{-35}$\,m & $H=1$ Hopfion \\
2 & Baryons          & $10^{-15}$\,m & 3-Hopfion bound state \\
3 & Atoms            & $10^{-10}$\,m & OHC Bessel modes \\
4 & Molecules        & $10^{-9}$\,m  & Hopf channel overlap \\
5 & Macromolecules   & $10^{-6}$\,m  & High-$H$ Hopf structures \\
6 & Stars            & $10^{8}$\,m   & Hopf pressure balance \\
7 & Galaxies         & $10^{21}$\,m  & OHC vortex topology \\
8 & Cosmic web       & $10^{24}$\,m  & Cosmological vortex network \\
9 & Universe         & $10^{26}$\,m  & Closed OHC multiplet ($H=0$) \\
C & Consciousness    & $10^{-1}$\,m  & $\mathcal{H} > S_{D_4} n_\mathrm{BCT}$ \\
\bottomrule
\end{tabular}
\end{table}

Scale C is not an anomaly in this table. It is the
inevitable consequence of the same theorem applied to the
one class of physical structure that can ask the theorem about
itself.

\section{Discussion}

Three points warrant explicit statement.

First: the cosmic web scale prediction
$\lambda_\mathrm{web} \approx 87$\,Mpc
from Eq.~(\ref{eq:web}) is a parameter-free BCT prediction
testable against large-scale structure surveys (DESI, Euclid,
Vera Rubin Observatory). The BCT prediction uses only the dark
energy scale derived in Letter~5; no additional parameters are
introduced.

Second: the claim about consciousness in
Eq.~(\ref{eq:crit}) is a threshold condition, not a theory
of qualia or phenomenal experience. The BCT programme claims
only that there is a \textit{physical} threshold --- derivable
from vacuum geometry --- that separates systems capable of
topological self-modelling from those that are not.
Whether this threshold is sufficient for consciousness is a
philosophical question; whether it is necessary is an
empirical one.

Third: the total Hopf charge $H_\mathrm{tot} = 0$ for the
universe is a statement about the observable universe as
a physical system, not a claim about what lies beyond the
de Sitter horizon. The BCT vacuum is infinite; the observable
universe is one OHC multiplet within it.

\section{Conclusion}

The BCT Persistence Theorem, established in Letter~69 for
scales 1 through 6, is here extended to its maximum scope.
Every persistent structure in the physical universe ---
from a single Hopfion to the cosmic web to a conscious mind
--- is a topological enclosure of information by OHC vacuum
geometry.

The unifying equation is $H \in \mathbb{Z}$, the topological
quantisation of the Hopf fibration $\pi_3(S^2) = \mathbb{Z}$.

What exists is what the OHC can contain.

\begin{thebibliography}{99}
\bibitem{L69} M.~R.~Cabri\'{e}, BCT Letter~69: The Persistence
Theorem, Zenodo (2026).
\bibitem{L34} M.~R.~Cabri\'{e}, BCT Letter~34: All Three
Standard Model Scales, Zenodo DOI: 10.5281/zenodo.18957202 (2026).
\bibitem{L5}  M.~R.~Cabri\'{e}, BCT Letter~5: The Cosmological
Constant, Zenodo (2026).
\bibitem{L36} M.~R.~Cabri\'{e}, BCT Letter~36: The OHC
Interior, Zenodo DOI: 10.5281/zenodo.18974754 (2026).
\bibitem{L78} M.~R.~Cabri\'{e}, BCT Letters~78--85: Chemistry
Series, Zenodo (2026).
\bibitem{monograph} M.~R.~Cabri\'{e}, The BCT Superfluid
Lattice Model (Monograph), Zenodo
DOI: 10.5281/zenodo.18884976 (2026).
\end{thebibliography}

\end{document}
